On considère le pavé droit ABCDEFGH tel que AB $= 3$ et AD = AE $= 1$ représenté ci-dessous.

\begin{center}
\psset{unit=1cm}
\begin{pspicture}(9.4,6.3)
\psframe(0.3,0.6)(7.8,4.3)
\psline(7.8,0.6)(8.6,2.2)(8.6,5.9)(7.8,4.3)
\psline(8.6,5.9)(1.1,5.9)(0.3,4.3)
\psline[linestyle=dashed](0.3,0.6)(1.1,2.2)(8.6,2.2)
\psline[linestyle=dashed](1.1,2.2)(1.1,5.9)
\psdots[dotstyle=+,dotangle=45,dotscale=1.9](2.8,0.6)(4.85,2.2)\uput[d](2.8,0.6){I}
\uput[dl](0.3,0.6){A} \uput[dr](7.8,0.6){B} \uput[r](8.6,2.2){C} \uput[ul](1.1,2.2){D}
\uput[l](0.3,4.3){E} \uput[r](7.8,4.3){F} \uput[ur](8.6,5.9){G} \uput[ul](1.1,5.9){H}
\uput[u](4.85,2.2){M}
\end{pspicture}
\end{center}

On considère le point I du segment [AB] tel que $\vect{\text{AB}}= 3 \vect{\text{AI}}$ et on appelle $M$ le milieu du segment [CD].

On se place dans le repère orthonormé $\left(\text{A}~;~ \vect{\text{AI}}, \vect{\text{AD}}, \vect{\text{AE}}\right)$.

\medskip

\begin{enumerate}
\item Sans justifier, donner les coordonnées des points F{}, H et M.

\item 
	\begin{enumerate}
		\item 	Montrer que le vecteur $\vect{n}\begin{pmatrix}2\\6\\ 3\end{pmatrix}$ est un vecteur normal au plan (HMF).
		\item En déduire qu'une équation cartésienne du plan (HMF) est :

\[2x + 6y + 3z - 9 = 0.\]

		\item Le plan $\mathcal{P}$ dont une équation cartésienne est $5x + 15y - 3z + 7 = 0$ est-il parallèle au plan (HMF) ? Justifier la réponse.
	\end{enumerate}
\item Déterminer une représentation paramétrique de la droite (DG).
\item On appelle $N$ le point d'intersection de la droite (DG) avec le plan (HMF).

Déterminer les coordonnées du point N.
\item Le point R de coordonnées $\left(3~;~\dfrac{1}{4}~;~\dfrac{1}{2}\right)$ est-il le projeté orthogonal du point G sur le plan (HMF) ? Justifier la réponse.
\end{enumerate}
